

\newpage

\paragraph{Problem 1 answer}\GradescopeComment{This should be on page 3 (top) of your PDF.}

%%%%%%%%%%%%%%%%%%%%%%%%%%%%%%%%%%%%%%%%%%%%%%%%%%%%%% 
%%%%%%%%%%%%%%%%%%%%%%%%%%%%%%%%%%%%%%%%%%%%%%%%%%%%%% 
%% 
%% PROBLEM 1(a) ANSWER 
%% 
%%%%%%%%%%%%%%%%%%%%%%%%%%%%%%%%%%%%%%%%%%%%%%%%%%%%%% 
%%%%%%%%%%%%%%%%%%%%%%%%%%%%%%%%%%%%%%%%%%%%%%%%%%%%%% 

\begin{answerBox}{3.8in}   % please don't edit this line, doing so will throw off the page layout 

  \paragraph{(a)}  

  % Define your reduction $r$
  % by defining, for every possible instance $I=(p, q, a, b)$ of  \prob{Family Seating},
  % what output $(G, s, t, k) = r(I)$ the reduction $r$ gives.
  % Your definition must state precisely how $G$, $s$, $t$, and $k$ are defined for the given $(p, q, a, b)$.

  Fix an instance $I=(p, q, a=(a_1,a_2,\ldots,a_p), b=(b_1,b_2,\ldots,b_q))$ of  \prob{Family Seating}.  
  The output of the reduction $r$ on input $I$ is $r(I) = (G=(V, E), s, t, k)$ where
  digraph $G=(V,E)$ with edge capacities, source $s$, sink $t$, and value $k$ are defined as follows.

  \bigskip 

  \begin{blue}

    % Explain as precisely as you can in words.
    % Use formal notation only when necessary for precision.

    \REPLACEME

  \end{blue}

\end{answerBox}

%%%%%%%%%%%%%%%%%%%%%%%%%%%%%%%%%%%%%%%%%%%%%%%%%%%%%% 
%%%%%%%%%%%%%%%%%%%%%%%%%%%%%%%%%%%%%%%%%%%%%%%%%%%%%% 
%% 
%% PROBLEM 1(b) ANSWER 
%% 
%%%%%%%%%%%%%%%%%%%%%%%%%%%%%%%%%%%%%%%%%%%%%%%%%%%%%% 
%%%%%%%%%%%%%%%%%%%%%%%%%%%%%%%%%%%%%%%%%%%%%%%%%%%%%% 

\begin{answerBox}{1.8in}   % please don't edit this line, doing so will throw off the page layout 

  \paragraph{(b)}  

  % Show that your reduction is a \emph{polynomial-time} reduction,
  % by explaining carefully how, given $I$, one can compute the output $r(I)$ in time polynomial in the input size ($p+q$).

  \begin{blue}

    \REPLACEME

  \end{blue}

\end{answerBox}

%%%%%%%%%%%%%%%%%%%%%%%%%%%%%%%%%%%%%%%%%%%%%%%%%%%%%% 
%%%%%%%%%%%%%%%%%%%%%%%%%%%%%%%%%%%%%%%%%%%%%%%%%%%%%% 
%% 
%% PROBLEM 1(c) ANSWER 
%% 
%%%%%%%%%%%%%%%%%%%%%%%%%%%%%%%%%%%%%%%%%%%%%%%%%%%%%% 
%%%%%%%%%%%%%%%%%%%%%%%%%%%%%%%%%%%%%%%%%%%%%%%%%%%%%% 

\begin{answerBox}{3in}   % please don't edit this line, doing so will throw off the page layout 

  \paragraph{(c)}

  % Consider the example input $I$ with $p=2$, $q=3$, $a=(2,2)$, $b=(1,2,2)$.
  % Given this specific $I$, calculate the output $(G, s, t, k) = r(I)$.
  % Draw the resulting flow network.  What is $k$? 

  For the input $I$ with $p=2$, $q=3$, $a=(2,2)$, $b=(1,2,2)$, the output $r(I)$ of the reduction is as follows.
  The value of $k$ is
  \BLUE{
    \REPLACEME
    .}

  Here's a drawing of the flow network $G$:

  \bigskip
  
  \begin{center}

    \begin{blue}

      \REPLACEME
      
      % % Here's an example illustrating one way to render your graph in LaTeX

      % \begin{tikzpicture}[
      %   vertex/.style={circle, inner sep=0pt, minimum width=15pt, draw}, 
      %   label/.style={pos=0.3, fill=white, inner sep=2pt}, 
      %   every path/.style={->}
      %   ]

      %     %%   NODES 
      
      %   \node[vertex] (s) {s}; 
      %   \node[vertex] (u1) [right =of s, yshift=20pt] {$u_1$}; 
      %   \node[vertex] (w1) [right =of u1] {$w_1$}; 
      %   \node[vertex] (t) [right =of w1, yshift=-20pt] {$t$}; 

      %     %%   EDGES AND CAPACITIES

      %   \path (s) edge node [label] {1} (u1); 
      %   \path (u1) edge node [label] {2}  (w1); 
      %   \path (w1) edge node [label] {3}  (t); 

      % \end{tikzpicture}

      % Or you can delete the block above, and instead draw the forest by hand,
      % scan it to a PDF, then include the PDF in your LaTeX document
      % using these commands:

      % \includegraphics[height=2.5in,width=6in]{YOUR_FILE_NAME_HERE}   % FILE NAME WITHOUT THE PDF EXTENSION

      % See https://texblog.org/2012/02/23/crop-figures-with-includegraphics
      % if you want to crop the included image.

    \end{blue}
  \end{center}

\end{answerBox}

%%%%%%%%%%%%%%%%%%%%%%%%%%%%%%%%%%%%%%%%%%%%%%%%%%%%%% 
%%%%%%%%%%%%%%%%%%%%%%%%%%%%%%%%%%%%%%%%%%%%%%%%%%%%%% 
%% 
%% PROBLEM 1(d) ANSWER 
%% 
%%%%%%%%%%%%%%%%%%%%%%%%%%%%%%%%%%%%%%%%%%%%%%%%%%%%%% 
%%%%%%%%%%%%%%%%%%%%%%%%%%%%%%%%%%%%%%%%%%%%%%%%%%%%%% 

\newpage

\paragraph{Problem 1 answer, continued}\GradescopeComment{This should be on page 4 (top) of your PDF.}

\begin{answerBox}{3in}   % please don't edit this line, doing so will throw off the page layout 
  
  \paragraph{(d)}

  % Show the ``if'' direction.
  % Given an arbitrary \prob{Family Seating} instance $I$,
  % let $G=(V,E)$, $s, t$, and $k$ be the output of your reduction $r$ given $I$.
  % Describe carefully how, given any $s$-$t$ flow $f$ in $G$ of value at least $k$,
  % one can obtain a valid seating arrangement for $I$.
  % (Assume that, if all edge capacities in $G$ are integers, then on each edge $(u,w)\in E$ the flow $f(u,w)$ is an integer.)

  Fix an arbitrary \prob{Family Seating} instance $I=(p, q, a=(a_1, \ldots, a_p), b=(b_1,\ldots, b_q))$.
  Let $G=(V,E)$, $s, t$, and $k$ be the output $r(I)$ of the reduction $r$ given $I$.
  Assume that $G$ has an integer $s$-$t$ flow of value at least $k$.
  Let $f$ be such a flow.
  From $f$, a valid seating arrangement $X$ for $I$ can be obtained as follows:

  \bigskip 

  For each family $i\in\{1,\ldots,p\}$ and table $j\in\{1,\ldots,q\}$,
  make the number of members of family $i$ sitting at table $j$ be
  \BLUE{
    \REPLACEME
  .}

\end{answerBox}

%%%%%%%%%%%%%%%%%%%%%%%%%%%%%%%%%%%%%%%%%%%%%%%%%%%%%% 
%%%%%%%%%%%%%%%%%%%%%%%%%%%%%%%%%%%%%%%%%%%%%%%%%%%%%% 
%% 
%% PROBLEM 1(e) ANSWER 
%% 
%%%%%%%%%%%%%%%%%%%%%%%%%%%%%%%%%%%%%%%%%%%%%%%%%%%%%% 
%%%%%%%%%%%%%%%%%%%%%%%%%%%%%%%%%%%%%%%%%%%%%%%%%%%%%% 

\begin{answerBox}{5in}   % please don't edit this line, doing so will throw off the page layout 

  \paragraph{(e)}

  % 
  % For the example input $I$ from Part B, the answer is \textsf{yes},
  % so the flow network that you drew in Part B should have a flow of value at least $k$ (for that $k$).
  % What is a flow in that flow network of value at least $k$?
  % What is a corresponding seating arrangement for the example input $I$?
  % 
  For the flow network in Part (c), a flow of value at least $k$ is as follows:

  The value of $k$ is
  \BLUE{
    \(
    \REPLACEME
    \).
  }

  Here is a drawing of a flow of value $k$ in that network:
  \medskip

  \begin{center}
    \begin{blue}

      \REPLACEME

      % % Here's an example illustrating one way to render your graph in LaTeX

      % \begin{tikzpicture}[
      %   vertex/.style={circle, inner sep=0pt, minimum width=15pt, draw}, 
      %   label/.style={pos=0.3, fill=white, inner sep=2pt}, 
      %   every path/.style={->}
      %   ]

      %     %%   NODES 
      
      %   \node[vertex] (s) {s}; 
      %   \node[vertex] (u1) [right =of s, yshift=20pt] {$u_1$}; 
      %   \node[vertex] (w1) [right =of u1] {$w_1$}; 
      %   \node[vertex] (t) [right =of w1, yshift=-20pt] {$t$}; 

      %     %%   EDGES AND FLOWS

      %   \path (s) edge node [label] {1} (u1); 
      %   \path (u1) edge node [label] {2}  (w1); 
      %   \path (w1) edge node [label] {3}  (t); 

      % \end{tikzpicture}

      % Or you can delete the block above, and instead draw the forest by hand,
      % scan it to a PDF, then include the PDF in your LaTeX document
      % using these commands:

      % \includegraphics[height=2.5in,width=6in]{YOUR_FILE_NAME_HERE}   % FILE NAME WITHOUT THE PDF EXTENSION

      % See https://texblog.org/2012/02/23/crop-figures-with-includegraphics
      % if you want to crop the included image.

    \end{blue}
  \end{center}

  % What is a corresponding seating arrangement for the example input $I$?

  \medskip

  \BLUE{A corresponding seating arrangement is as follows:}

  \begin{blue}

    \REPLACEME

  \end{blue}

\end{answerBox}

%%%%%%%%%%%%%%%%%%%%%%%%%%%%%%%%%%%%%%%%%%%%%%%%%%%%%% 
%%%%%%%%%%%%%%%%%%%%%%%%%%%%%%%%%%%%%%%%%%%%%%%%%%%%%% 
%% 
%% PROBLEM 1(f) ANSWER 
%% 
%%%%%%%%%%%%%%%%%%%%%%%%%%%%%%%%%%%%%%%%%%%%%%%%%%%%%% 
%%%%%%%%%%%%%%%%%%%%%%%%%%%%%%%%%%%%%%%%%%%%%%%%%%%%%% 

\newpage

\paragraph{Problem 1 answer, continued}\GradescopeComment{This should be on page 5 (top) of your PDF.}

\paragraph{(f)} 

% Now show the ``only if'' direction.
% Given an arbitrary \prob{Family Seating} instance $I$,
% let $G=(V,E)$, $s, t$, and $k$ be the output of $r$ given $I$.
% Describe carefully how, given any valid seating arrangement for $I$,
% one can obtain an $s$-$t$ flow $f$ in $G$ of value at least $k$.

Fix an arbitrary \prob{Family Seating} instance $I=(p, q, a=(a_1, \ldots, a_p), b=(b_1,\ldots, b_q))$.
Let $G=(V,E)$, $s, t$, and $k$ be the output $r(I)$ of the reduction $r$ given $I$.
Assume that $I$ has a valid seating arrangement.
Let $X$ be such an arrangement.
From $X$, an $s$-$t$ flow $f$ of value at least $k$ in $G$ can be obtained as follows.

\begin{blue}

  \REPLACEME

\end{blue}

%%%%%%%%%%%%%%%%%%%%%%%%%%%%%%%%%%%%%%%%%%%%%%%%%%%%%% 
%%%%%%%%%%%%%%%%%%%%%%%%%%%%%%%%%%%%%%%%%%%%%%%%%%%%%% 
%% 
%% PROBLEM 1 COLLABORATORS AND CITATIONS 
%% 
%%%%%%%%%%%%%%%%%%%%%%%%%%%%%%%%%%%%%%%%%%%%%%%%%%%%%% 
%%%%%%%%%%%%%%%%%%%%%%%%%%%%%%%%%%%%%%%%%%%%%%%%%%%%%% 

\vfill

\begin{answerBox}{2in}          % please don't edit this line, doing so will throw off the page layout
  
  \paragraph{Problem \arabic{problem} collaboration and citations}~\GradescopeComment
  {This should be on page 5 (bottom) of your PDF.}

  I collaborated on this problem with the following people:
  \begin{blue}
    \begin{itemize}
    \item \REPLACEME                    % write "none" if none
    \end{itemize}
  \end{blue}

  My answers use ideas from the following sources (other than the lecture notes and textbooks): 
  \begin{blue}
    \begin{itemize}
    \item \REPLACEME                    % write "none" if none
    \end{itemize}
  \end{blue}

\end{answerBox}

%%% Local Variables:
%%% mode: latex
%%% TeX-master: "homework_11"
%%% End:
